% GENERATED FILE. Edit canonical lessons, then run npm run generate:books.
% canonical-source-hash: fnv1a64:95c858414a91f54b
% canonical-lessons: GU-C266-lagni, GU-C266-sukh, GU-C266-darvu, GU-C266-radvu, GU-C266-sharmavu, GU-R266-review-261-266

\chapter{Fear, tears and shyness}
\label{ch:gu-fear-tears-and-shyness}
\hlchaptermodality{266}

\begin{chapteropening}
\textbf{By the end of this chapter:} \emph{I can say what someone is not afraid of and ask whether someone is crying.}

Complete the last lesson of chapter 266: I can say what someone is not afraid of and ask whether someone is crying.
\end{chapteropening}

\section[\texorpdfstring{\gu{લાગણી} (lāgṇī)}{લાગણી (lāgṇī)}]{\gu{લાગણી} (lāgṇī) — feeling, emotion}

\label{lesson:GU-C266-lagni}



\begin{quote}
Before the new one: say the Gujarati for daughter-in-law, then the Gujarati for trust.
\end{quote}

\subsection*{You'll want to know: \gu{લાગણી}}
\textbf{\gu{લાગણી}} — \emph{lāgṇī} — \textquotedblleft{}feeling, emotion\textquotedblright{}.

\begin{quote}
\textbf{\gu{આ} \gu{લાગણી} \gu{સારી} \gu{નથી}.} — \emph{ā lāgṇī sārī nathī} — \textquotedblleft{}This feeling is not good.\textquotedblright{}
\end{quote}

\begin{cousinweb}
\textbf{\gu{લાગણી}} (\emph{lāgṇī}), what touches you, is built on the inherited \emph{lāgvũ}, \textquotedblleft{}to come into contact with\textquotedblright{}, with \emph{-aṇī}. CDIAL derives \emph{lāgvũ} from Sanskrit \emph{lagyati}, \textquotedblleft{}adheres to\textquotedblright{} (CDIAL 10895).
\end{cousinweb}

\begin{grammarlens}[title={What you've built}]
One more word for everyday life at A2.
\end{grammarlens}

\subsection*{Your turn}
\begin{itemize}
\raggedright
  \item \emph{Say it:} \emph{lāgṇī}
  \item \emph{Say it:} \emph{lāgṇī}, once more
  \item \emph{Say it:} \emph{ā lāgṇī sārī nathī} — This feeling is not good.
  \item \emph{Recall:} say the Gujarati for daughter-in-law, then the Gujarati for trust, then say \emph{lāgṇī} again
\end{itemize}

\begin{tcolorbox}[breakable,colback=teal!4,colframe=teal!35!black,title={Before you move on}]
What does \gu{લાગણી} mean? (\textquotedblleft{}Feeling, emotion\textquotedblright{}.) Say it once more.
\end{tcolorbox}
\section[\texorpdfstring{\gu{સુખ} (sukh)}{સુખ (sukh)}]{\gu{સુખ} (sukh) — contentment, happiness}

\label{lesson:GU-C266-sukh}



\begin{quote}
Before the new one: say the Gujarati for trust, then the Gujarati for feeling.
\end{quote}

\subsection*{You'll want to know: \gu{સુખ}}
\textbf{\gu{સુખ}} — \emph{sukh} — \textquotedblleft{}contentment, happiness\textquotedblright{}.

\begin{quote}
\textbf{\gu{શું} \gu{પૈસામાં} \gu{સુખ} \gu{છે}?} — \emph{śũ paisāmā̃ sukh che?} — \textquotedblleft{}Is there happiness in money?\textquotedblright{}
\end{quote}

A neuter noun.

\begin{cousinweb}
\textbf{\gu{સુખ}} (\emph{sukh}) is Sanskrit \emph{sukhá}, \textquotedblleft{}pleasant; comfort, happiness\textquotedblright{}, taken over as a tatsama (CDIAL 13451). The same entry lists the inherited \emph{sahelũ}, \textquotedblleft{}easy\textquotedblright{}.
\end{cousinweb}

\begin{grammarlens}[title={What you've built}]
One more word for everyday life at A2.
\end{grammarlens}

\subsection*{Your turn}
\begin{itemize}
\raggedright
  \item \emph{Say it:} \emph{sukh}
  \item \emph{Say it:} \emph{sukh}, once more
  \item \emph{Say it:} \emph{śũ paisāmā̃ sukh che?} — Is there happiness in money?
  \item \emph{Recall:} say the Gujarati for trust, then the Gujarati for feeling, then say \emph{sukh} again
\end{itemize}

\begin{tcolorbox}[breakable,colback=teal!4,colframe=teal!35!black,title={Before you move on}]
What does \gu{સુખ} mean? (\textquotedblleft{}Contentment, happiness\textquotedblright{}.) Say it once more.
\end{tcolorbox}
\section[\texorpdfstring{\gu{ડરવું} (ḍarvũ)}{ડરવું (ḍarvũ)}]{\gu{ડરવું} (ḍarvũ) — to be afraid, to fear}

\label{lesson:GU-C266-darvu}



\begin{quote}
Before the new one: say the Gujarati for feeling, then the Gujarati for contentment.
\end{quote}

\subsection*{You'll want to know: \gu{ડરવું}}
\textbf{\gu{ડરવું}} — \emph{ḍarvũ} — \textquotedblleft{}to be afraid, to fear\textquotedblright{}.

\begin{quote}
\textbf{\gu{હું} \gu{અંધારાથી} \gu{નથી} \gu{ડરતો}.} — \emph{hũ andhārāthī nathī ḍarto} — \textquotedblleft{}I am not afraid of the dark.\textquotedblright{}
\end{quote}

\begin{cousinweb}
\textbf{\gu{ડરવું}} (\emph{ḍarvũ}) is a tadbhava. It comes from Sanskrit \emph{darati}, \textquotedblleft{}is frightened\textquotedblright{}; CDIAL glosses the Gujarati verb \textquotedblleft{}to fear\textquotedblright{} (CDIAL 6190).
\end{cousinweb}

\begin{grammarlens}[title={What you've built}]
One more word for everyday life at A2.
\end{grammarlens}

\subsection*{Your turn}
\begin{itemize}
\raggedright
  \item \emph{Say it:} \emph{ḍarvũ}
  \item \emph{Say it:} \emph{ḍarvũ}, once more
  \item \emph{Say it:} \emph{hũ andhārāthī nathī ḍarto} — I am not afraid of the dark.
  \item \emph{Recall:} say the Gujarati for feeling, then the Gujarati for contentment, then say \emph{ḍarvũ} again
\end{itemize}

\begin{tcolorbox}[breakable,colback=teal!4,colframe=teal!35!black,title={Before you move on}]
What does \gu{ડરવું} mean? (\textquotedblleft{}To be afraid, to fear\textquotedblright{}.) Say it once more.
\end{tcolorbox}
\section[\texorpdfstring{\gu{રડવું} (raḍvũ)}{રડવું (raḍvũ)}]{\gu{રડવું} (raḍvũ) — to cry, to weep}

\label{lesson:GU-C266-radvu}



\begin{quote}
Before the new one: say the Gujarati for contentment, then the Gujarati for to be afraid.
\end{quote}

\subsection*{You'll want to know: \gu{રડવું}}
\textbf{\gu{રડવું}} — \emph{raḍvũ} — \textquotedblleft{}to cry, to weep\textquotedblright{}.

\begin{quote}
\textbf{\gu{શું} \gu{બાળક} \gu{રડે} \gu{છે}?} — \emph{śũ bāḷak raḍe che?} — \textquotedblleft{}Is the baby crying?\textquotedblright{}
\end{quote}

\begin{cousinweb}
\textbf{\gu{રડવું}} (\emph{raḍvũ}) is a tadbhava. It comes from Sanskrit \emph{raṭati}, \textquotedblleft{}cries, yells\textquotedblright{}; CDIAL lists Gujarati \emph{raṛvũ}, \textquotedblleft{}to weep\textquotedblright{} (CDIAL 10590).
\end{cousinweb}

\begin{grammarlens}[title={What you've built}]
One more word for everyday life at A2.
\end{grammarlens}

\subsection*{Your turn}
\begin{itemize}
\raggedright
  \item \emph{Say it:} \emph{raḍvũ}
  \item \emph{Say it:} \emph{raḍvũ}, once more
  \item \emph{Say it:} \emph{śũ bāḷak raḍe che?} — Is the baby crying?
  \item \emph{Recall:} say the Gujarati for contentment, then the Gujarati for to be afraid, then say \emph{raḍvũ} again
\end{itemize}

\begin{tcolorbox}[breakable,colback=teal!4,colframe=teal!35!black,title={Before you move on}]
What does \gu{રડવું} mean? (\textquotedblleft{}To cry, to weep\textquotedblright{}.) Say it once more.
\end{tcolorbox}
\section[\texorpdfstring{\gu{શરમાવું} (śarmāvũ)}{શરમાવું (śarmāvũ)}]{\gu{શરમાવું} (śarmāvũ) — to feel shy, to be embarrassed}

\label{lesson:GU-C266-sharmavu}



\begin{quote}
Before the new one: say the Gujarati for to be afraid, then the Gujarati for to cry.
\end{quote}

\subsection*{You'll want to know: \gu{શરમાવું}}
\textbf{\gu{શરમાવું}} — \emph{śarmāvũ} — \textquotedblleft{}to feel shy, to be embarrassed\textquotedblright{}.

\begin{quote}
\textbf{\gu{નાની} \gu{છોકરી} \gu{શરમાતી} \gu{નથી}.} — \emph{nānī chokrī śarmātī nathī} — \textquotedblleft{}The little girl is not shy.\textquotedblright{}
\end{quote}

\begin{cousinweb}
\textbf{\gu{શરમાવું}} (\emph{śarmāvũ}) is built on \emph{śaram}, \textquotedblleft{}shame, shyness\textquotedblright{}, with \emph{-ā-}. \emph{Śaram} is a loan usually traced to Persian \emph{sharm}, \textquotedblleft{}shame\textquotedblright{}; CDIAL has no entry for it.
\end{cousinweb}

\begin{grammarlens}[title={What you've built}]
One more word for everyday life at A2.
\end{grammarlens}

\subsection*{Your turn}
\begin{itemize}
\raggedright
  \item \emph{Say it:} \emph{śarmāvũ}
  \item \emph{Say it:} \emph{śarmāvũ}, once more
  \item \emph{Say it:} \emph{nānī chokrī śarmātī nathī} — The little girl is not shy.
  \item \emph{Recall:} say the Gujarati for to be afraid, then the Gujarati for to cry, then say \emph{śarmāvũ} again
\end{itemize}

\begin{tcolorbox}[breakable,colback=teal!4,colframe=teal!35!black,title={Before you move on}]
What does \gu{શરમાવું} mean? (\textquotedblleft{}To feel shy, to be embarrassed\textquotedblright{}.) Say it once more.
\end{tcolorbox}
\section[(dialogue)]{Words from chapters 261–266 — a review}

\label{lesson:GU-R266-review-261-266}



\begin{quote}
Say the Gujarati for kite and to fly a kite.
\end{quote}

\subsection*{Your turn}
Say these aloud:

\begin{itemize}
\raggedright
  \item kite, to fly a kite, fair, garba, firecracker
  \item to hum, to call out, sports, picnic, to dive
  \item sister-in-law, maternal aunt, paternal aunt, mother-in-law, nephew
  \item neighbour, pal, female friend, to recognize, to embrace
  \item to love, to marry, engagement, daughter-in-law, trust
  \item feeling, contentment, to be afraid, to cry, to feel shy
\end{itemize}

\begin{tcolorbox}[breakable,colback=teal!4,colframe=teal!35!black,title={Before you move on}]
Kite? (\textbf{\gu{પતંગ}}.) To fly a kite? (\textbf{\gu{ઉડાડવું}}.) Fair? (\textbf{\gu{મેળો}}.) Garba? (\textbf{\gu{ગરબો}}.) Firecracker? (\textbf{\gu{ફટાકડો}}.) To hum? (\textbf{\gu{ગણગણવું}}.) To call out? (\textbf{\gu{પોકારવું}}.) Sports? (\textbf{\gu{રમતગમત}}.) Picnic? (\textbf{\gu{ઉજાણી}}.) To dive? (\textbf{\gu{ડૂબકી} \gu{મારવી}}.) Sister-in-law? (\textbf{\gu{ભાભી}}.) Maternal aunt? (\textbf{\gu{માસી}}.) Paternal aunt? (\textbf{\gu{ફોઈ}}.) Mother-in-law? (\textbf{\gu{સાસુ}}.) Nephew? (\textbf{\gu{ભત્રીજો}}.) Neighbour? (\textbf{\gu{પડોશી}}.) Pal? (\textbf{\gu{દોસ્ત}}.) Female friend? (\textbf{\gu{બહેનપણી}}.) To recognize? (\textbf{\gu{ઓળખવું}}.) To embrace? (\textbf{\gu{ભેટવું}}.) To love? (\textbf{\gu{પ્રેમ} \gu{કરવો}}.) To marry? (\textbf{\gu{પરણવું}}.) Engagement? (\textbf{\gu{સગાઈ}}.) Daughter-in-law? (\textbf{\gu{વહુ}}.) Trust? (\textbf{\gu{વિશ્વાસ}}.) Feeling? (\textbf{\gu{લાગણી}}.) Contentment? (\textbf{\gu{સુખ}}.) To be afraid? (\textbf{\gu{ડરવું}}.) To cry? (\textbf{\gu{રડવું}}.) To feel shy? (\textbf{\gu{શરમાવું}}.)
\end{tcolorbox}
